%%%%%%%%%%%%%%%%%%%%%%%%%%%%%%%%%%%%%%%%%%%%%%%%%%%%%%%%%%%%%
\section{模型的改进与推广}

\subsection{模型的改进方向}

\textbf{折射率包络自反演路线。} 当前 $n_1$ 依赖外部文献值，在文献缺失场景受限。利用反射率谱的上下包络同时反演 $n(\omega)$ 和 $d$ 可摆脱文献依赖。但该方法要求足够多的条纹数且对背景扣除方案敏感，单晶圆数据可能仍面临 $n$-$d$ 简并，需变角测量系列破解。

\textbf{Lorentz 振子色散模型。} 采用 Drude-Lorentz 模型显式建模剩余射线带的两振子（793.9 和 970.1\,cm$^{-1}$），物理完整性高于纯 Drude 模型，有望收窄当前 $\sigma_p$ 的双口径张力，从而收窄厚度系统区间。

\textbf{过渡层与粗糙度显式建模。} v2.0 Si 拟合残差仍含 8\%--10\% 结构，可能是掺杂过渡层或界面微观粗糙度所致。在 Airy 式(12)中引入等效界面过渡矩阵可覆盖这一效应，但需警惕额外自由参数导致的可辨识性下降。

\subsection{模型的推广}

\textbf{多层结构。} Airy 三介质公式(12)可推广至传输矩阵法（Transfer-Matrix Method, TMM），处理任意多层膜系。这在 SOI、SiC/Si 异质结等实际器件结构中具有直接应用价值。

\textbf{变角测量系列。} 多角度（10\degree 至 70\degree，步长 5\degree）的反射率谱可提供 $n(\sigma)$ 与 $d$ 的独立辨识信息——入射角变化使 $\cos\theta_1$ 充分变化，$n$-$d$ 乘积不变性被打破，可实现文献无依赖的自洽反演。

\textbf{其他半导体材料。} 本文的方法框架（三域互证 + 两级一致性检验 + 多光束效应分离）不限于 SiC，可直接推广至 Si、GaAs、GaN 等外延材料的厚度测量，前提是适用条件（透明层与光学厚衬底）满足。模型参数需替换为相应材料的文献值。